% =====================================================================
% DDLab sense-extraction benchmark — workshop paper, first draft.
%
% STATUS: draft. Compiles with pdflatex + bibtex (standard article class;
% retarget to a specific workshop style file before submission).
%
% GAPS TO CLOSE BEFORE SUBMISSION:
%  1. Related work — verify every reference in references.bib (years,
%     venues, authors) against sources; expand the positioning vs
%     Othello-GPT / emergent world representations and bisimulation work.
%  2. Section 4 (Z-spectrum) — the edge-of-chaos lift is partly confounded
%     with label degeneracy at the ordered/chaotic extremes; it is framed
%     here as a secondary observation, not a headline. Keep it that way or
%     cut it.
%  3. Figures — three placeholders (\figtodo) to replace with real plots:
%     basin-field visualization, AMI-vs-K cliff, perceptual-mismatch bars.
%  4. Optional: add a sixth, non-contrastive SSL objective (BYOL/VICReg)
%     to strengthen the "structural, not optimization" claim.
%  5. Scale (n<=14) is a deliberate design choice (exact ground truth
%     requires enumerability) — stated in Limitations; keep owning it.
% =====================================================================
\documentclass[11pt]{article}

\usepackage[utf8]{inputenc}
\usepackage[T1]{fontenc}
\usepackage[margin=1in]{geometry}
\usepackage{amsmath,amssymb}
\usepackage{graphicx}
\usepackage{booktabs}
\usepackage{xcolor}
\usepackage[numbers,sort&compress]{natbib}
\usepackage[hidelinks]{hyperref}

\newcommand{\figtodo}[1]{%
  \begin{center}\fbox{\parbox{0.86\linewidth}{\centering\vspace{1.6em}%
  \textit{[FIGURE TODO: #1]}\vspace{1.6em}}}\end{center}}
\newcommand{\succf}{\mathrm{succ}}

\title{Perception Is Not Causation:\\
Auditing Representation Learners Against Exact Causal Ground Truth}
\author{Zaccarie Aglas\\\texttt{zaccarie@zaccarie.fr}}
\date{Draft --- \today}

\begin{document}
\maketitle

\begin{abstract}
Representation learning is usually evaluated without access to the causal
structure of the data: a probe reaches 85\% accuracy and we cannot tell
whether a feature is genuinely represented or the number is an artifact. We
introduce a testbed in which the causal structure is known \emph{exactly}. On
elementary cellular automata and random Boolean networks small enough to
enumerate, the entire dynamics --- basins of attraction, attractors,
garden-of-Eden states, the bisimulation quotient --- is computed exactly and
used as ground truth. With it we show three things. (i) A model trained only
to predict the forward dynamics develops a representation that linearly
encodes \emph{local} causal facts (garden-of-Eden membership, in-degree).
(ii) A compressed, structure-preserving discretization of the state space
provably \emph{exists} --- the dynamical quotient. (iii) Yet five standard
learned-tokenizer objectives --- reconstruction, next-step prediction,
multi-horizon prediction, and two temporal-contrastive variants --- all fail
to recover it. We explain the failure: the dynamical quotient is
\emph{perceptually shattered}, as dispersed in observation space as a random
partition, so objectives biased toward observation proximity optimize a
geometry orthogonal to the causal one. The barrier is structural, not an
optimization artifact. A single inductive bias --- locality --- unifies the
positive and negative results.
\end{abstract}

\section{Introduction}

Probing studies ask whether a learned representation encodes some property of
the data. They report an accuracy. But without the true structure of the
data, that accuracy is hard to calibrate: it conflates ``the feature is
explicitly represented'' with ``a powerful probe can compute the feature'' and
with chance-level artifacts. Studies of emergent world representations face
the same gap --- the ground truth (e.g.\ a board state) is known, but the
\emph{dynamical} structure (which states are causally equivalent, which are
unreachable, how trajectories converge) is not enumerated.

We remove the gap. We study discrete dynamical systems small enough to
enumerate completely. For an elementary cellular automaton on $n$ cells there
are $2^n$ states; for $n \le 20$ we build the entire state-transition graph
and from it compute, \emph{exactly and without sampling}, every dynamical
object of interest: attractors, basins of attraction, pre-images,
garden-of-Eden states, transient lengths, and the bisimulation quotient. This
exact ground truth lets us audit a representation learner the way one audits a
model against a known answer.

Using the testbed we establish three results.

\begin{itemize}
  \item \textbf{Forward training recovers local causal structure}
    (Section~\ref{sec:local}). A model trained only to predict the next state
    develops a representation from which \emph{local} causal facts ---
    garden-of-Eden membership, in-degree --- are linearly decodable, well
    above raw-observation and untrained-network baselines. Long-horizon
    structure (basin membership) is \emph{not} recovered.
  \item \textbf{The structure-preserving discretization exists}
    (Section~\ref{sec:exists}). A compressed tokenization that exactly
    preserves the basin structure provably exists: the dynamical quotient.
    We exhibit it constructively.
  \item \textbf{Learned tokenizers fail to recover it}
    (Section~\ref{sec:fail}). Five standard learned-tokenizer objectives all
    fail. Section~\ref{sec:why} explains why: the dynamical quotient is
    \emph{perceptually shattered}, so observation-proximity objectives are
    structurally mismatched to it.
\end{itemize}

A single mechanism --- an inductive bias toward locality in
observation space --- accounts for both the positive result (local structure
is recovered) and the negative one (the non-local quotient is missed).

\section{Related Work}
\label{sec:related}

\paragraph{Emergent world representations and probing.}
Probing for the latent state of a system underlying a sequence model is an
established paradigm; \citet{li2023emergent} train a sequence model on a
synthetic task and probe for the world state. Linear probes
\citep{alain2016probes} operationalize ``explicitly represented.'' We differ
in two ways: our ground truth is the full \emph{dynamical} structure (basins,
attractors, the bisimulation quotient), exactly enumerated; and our central
object of study is the systematic \emph{failure} of learned tokenizers, with a
structural explanation.

\paragraph{Bisimulation and state abstraction.}
Bisimulation metrics learn representations that collapse behaviorally
equivalent states \citep{zhang2021dbc,castro2021mico}. Our testbed gives a
setting in which the bisimulation quotient is known exactly, so such methods
can be validated against ground truth rather than evaluated only by downstream
return.

\paragraph{World models, tokenization, information.}
Discrete-latent and tokenized world models \citep{vandenoord2017vqvae,
hafner2023dreamerv3} motivate our ``learned tokenizer'' setting. The
Information Bottleneck \citep{tishby1999ib} and Shannon's separation of
transmission from meaning \citep{shannon1948} frame the compression questions.

\paragraph{Discrete dynamics.}
The exact analysis of cellular-automaton state space --- basins of attraction,
the $Z$ parameter, the edge of chaos --- follows \citet{wuensche1999z}.
Analysis of recurrent networks as dynamical systems
\citep{sussillo2013blackbox} is a methodological precedent.

% TODO: verify all citations; expand positioning vs Othello-GPT follow-ups
% and the contrastive / CPC literature \citep{vandenoord2018cpc}.

\section{The Testbed: Discrete Dynamics with Exact Ground Truth}
\label{sec:testbed}

\paragraph{Systems.}
An elementary cellular automaton (ECA) is a one-dimensional binary automaton
with a radius-1 neighborhood and periodic boundary; a rule is one of $256$
local update functions. A state of $n$ cells is an integer in $[0, 2^n)$. We
also use random Boolean networks (RBN) --- the Kauffman model, with irregular
wiring --- to check that results are not an artifact of the ECA's regular
topology.

\paragraph{Exact structure.}
For $n \le 20$ we enumerate all $2^n$ states and build the complete
state-transition graph $\succf : S \to S$. Because the graph is functional
(out-degree one), we compute exactly: the attractor cycles; the basin of
attraction of each state (the attractor it drains to); all pre-images;
garden-of-Eden states (no pre-image); transient lengths. The
\emph{dynamical quotient} is the partition of states by basin --- equivalently,
the bisimulation quotient under $\succf$. None of this is sampled.

\paragraph{Regimes.}
Rules span ordered, complex, and chaotic regimes, indexed by the $Z$ parameter
\citep{wuensche1999z}. We use rule~250 (ordered), rule~110 (complex), and
rule~30 (chaotic) as representatives.

\figtodo{Basin-of-attraction field of rule 110 ($n{=}9$): a force-directed
graph, nodes coloured by basin, attractor cycles as hubs, garden-of-Eden
states as leaves. Generated by the project's visualization script.}

\section{Forward Training Recovers Local Causal Structure}
\label{sec:local}

We train a small GRU to predict the next state from the current one --- the
forward (next-token) objective. It reaches one-step accuracy $1.0$ and
generalizes to held-out basins, so it has learned the rule rather than
memorized a table.

We then freeze it and probe its representation for causal facts that the
forward objective never required: garden-of-Eden membership and in-degree.
These are \emph{inverse} facts --- about a state's past, not its future.
Table~\ref{tab:probe} reports linear-probe test accuracy on rule~110, against
two baselines: the raw observation bits, and the same network untrained.

\begin{table}[h]
\centering
\caption{Linear-probe decoding of causal-inverse facts (rule 110, $n{=}9$).
Forward training makes garden-of-Eden membership linearly decodable well above
the raw-observation and untrained-network baselines.}
\label{tab:probe}
\begin{tabular}{lccc}
\toprule
target & raw bits & untrained net & forward-trained \\
\midrule
garden-of-Eden (acc.) & 0.68 & 0.68 & \textbf{0.86} \\
in-degree (acc.)      & 0.60 & 0.60 & \textbf{0.77} \\
\bottomrule
\end{tabular}
\end{table}

Forward training builds a representation in which \emph{local} causal facts
become linearly accessible. Two caveats keep this honest. First, an MLP probe
reaches $0.97$ on raw bits too: the facts are computable from the observation;
the forward representation makes them \emph{linearly} accessible, i.e.\
disentangled. Second, across the $Z$ spectrum the disentangling effect is
largest in the complex regime; we report this as a secondary observation, as
the comparison is partly confounded with label balance at the ordered and
chaotic extremes.

Crucially, the \emph{long-horizon} structure is \emph{not} recovered: the
clustering induced by the forward representation has near-zero adjusted mutual
information (AMI) with the true basin partition. Forward training recovers
local structure and misses global structure --- a distinction
Section~\ref{sec:why} explains.

\section{The Structure-Preserving Discretization Exists}
\label{sec:exists}

Does a tokenization exist that \emph{compresses} the state space yet
\emph{preserves} the basin structure? It does, and we exhibit it.

Define the horizon quotient $\phi_H(s) = \succf^{H}(s)$: two states share a
code iff they converge to the same state within $H$ steps. Converging states
share a basin, so $\phi_H$ \emph{never} merges two basins --- it is, for every
$H$, a refinement of the dynamical quotient. As $H$ grows, states reach their
attractor cycles and the number of distinct codes drops sharply.

\begin{table}[h]
\centering
\caption{The horizon quotient (rule 110, $n{=}9$, $512$ states; no model). It
compresses while AMI with the true basin partition stays exactly $1.0$.}
\label{tab:horizon}
\begin{tabular}{rccc}
\toprule
horizon $H$ & distinct codes & collision rate & AMI vs.\ true basins \\
\midrule
0  & 512 & 0.00 & 1.00 \\
2  & 200 & 0.61 & 1.00 \\
8  & 109 & 0.79 & 1.00 \\
16 &  73 & 0.86 & 1.00 \\
\bottomrule
\end{tabular}
\end{table}

A discretization that compresses $512$ states to $73$ while preserving the
basin structure \emph{exactly} therefore exists. The open question is not its
existence but whether it can be \emph{learned} from observations.

\section{Learned Tokenizers Fail to Recover It}
\label{sec:fail}

We render each state as an image and train a tokenizer on the images, then
push the true dynamics through the learned codes and measure how well the
induced latent transition graph preserves the basin structure (AMI), as a
function of codebook size $K$.

We test five objectives: image reconstruction; next-image prediction;
multi-horizon (far-future) image prediction; one-step temporal contrastive
(InfoNCE with the successor as positive); and far-future temporal contrastive,
designed to break local perceptual similarity. Table~\ref{tab:cliff} shows the
first two; the contrastive variants never exceed AMI $\approx 0.12$ at any $K$.

\begin{table}[h]
\centering
\caption{AMI vs.\ true basins for learned tokenizers (rule 110, $n{=}9$). At
$K{=}512$ the tokenization is near-lossless and structure is preserved; the
first forced compression ($K{=}256$) collapses it --- a cliff, not a slope.}
\label{tab:cliff}
\begin{tabular}{lcc}
\toprule
objective & AMI at $K{=}512$ & AMI at $K{=}256$ \\
\midrule
reconstruction       & 0.95 & 0.02 \\
next-image prediction & 0.99 & 0.07 \\
\bottomrule
\end{tabular}
\end{table}

No objective recovers the quotient under compression. All preserve structure
only when the tokenization is near-lossless; the first bit of forced
compression destroys it. We also encountered, and document, the training
pathologies of this regime --- codebook collapse, and an objective/anti-collapse
conflict --- which we control for; the cliff is not an artifact of them.

\section{Why: The Perceptual--Dynamical Mismatch}
\label{sec:why}

The five objectives share an inductive bias: they map perceptually similar
observations to similar codes. We measure whether that bias can possibly align
with the dynamical quotient.

For each state we take its perceptual neighbors --- the states one bit away ---
and measure the fraction that lie in a \emph{different} basin. We compare this
\emph{mismatch} to the value expected if basins were assigned at random with
the same sizes.

\begin{table}[h]
\centering
\caption{Perceptual--dynamical mismatch ($n{=}9$): fraction of one-bit
perceptual neighbors in a different basin, vs.\ a random-assignment baseline.
The ratio is $1.0$ in every regime.}
\label{tab:mismatch}
\begin{tabular}{lcccc}
\toprule
rule & regime & mismatch & random baseline & ratio \\
\midrule
250 & ordered  & 0.004 & 0.004 & 1.00 \\
110 & complex  & 0.914 & 0.902 & 1.01 \\
30  & chaotic  & 0.316 & 0.315 & 1.00 \\
\bottomrule
\end{tabular}
\end{table}

\figtodo{Perceptual--dynamical mismatch vs.\ the random baseline, bar chart per
regime --- the two bars are equal in every regime.}

The mismatch equals the random baseline. The basins are \emph{perceptually
shattered}: as dispersed in observation space as a random partition. Perceptual
proximity carries no information about basin membership beyond the base rate.

Hence the failure of Section~\ref{sec:fail} is \emph{structural}. An objective
biased toward observation proximity optimizes a geometry that is provably
orthogonal to the dynamical quotient --- it is not under-optimized, it is
optimized toward the wrong target. The image \emph{determines} the basin (it
determines the state); the barrier is not informational but inductive.

This unifies the paper. Section~\ref{sec:local} found that forward training
recovers \emph{local} causal structure; here we find the basin is perceptually
\emph{non-local}. One inductive bias --- locality --- explains both the local
success and the global failure.

\section{Discussion and Limitations}

\paragraph{Implications.}
Reconstruction and instance-contrastive latents should not be assumed to carry
causal-dynamical structure; where the target structure is non-local in
observation space, an objective with a locality bias is structurally
mismatched to it. This is an argument for objectives defined recursively on
the dynamics --- bisimulation-style --- over observation-proximity objectives.
The mismatch measure of Section~\ref{sec:why} is a cheap a-priori check: if a
target structure is perceptually shattered relative to the observation, a
perceptual objective will fail before any training is run. More broadly, the
testbed itself is a usable methodology --- audit a representation learner on a
small enumerable system where the causal answer is known.

\paragraph{Limitations.}
The systems are small ($n \le 14$ in experiments). This is deliberate: exact,
sampling-free ground truth requires an enumerable state space. The testbed is
a controlled audit, not a claim about any specific large model; transfer of
the mechanism to large, continuous, learned-latent world models is the natural
next question and is not established here. Several results use single
instances; the central negative (the perceptual-shatter ratio) is robust
across regimes, but ensemble averaging would strengthen the secondary
observations.

\section{Conclusion}

On discrete dynamical systems with exact ground truth, a compressed
structure-preserving discretization of the state space provably exists --- the
dynamical quotient --- and five standard representation-learning objectives
provably fail to recover it from observations. The reason is structural: the
quotient is perceptually shattered, and objectives biased toward observation
proximity optimize an orthogonal geometry. Perception is not causation, and an
objective that pursues the first will not find the second.

\paragraph{Reproducibility.}
All experiments are exact or seeded; the implementation (the testbed, the
eleven experiment phases, and a test suite) is released.

\bibliographystyle{plainnat}
\bibliography{references}

\end{document}
